\chapter{Nail Injuries}

\section*{Definition}
Nail injuries include a range of traumas to the fingernails or toenails, such as subungual hematoma (blood under the nail), nail avulsion (partial or complete detachment), nail bed laceration, and fractures of the fingertip involving the nail. These are common hand injuries seen in emergency departments.

\section*{What Happens in Your Body}
The nail is produced by the nail matrix, located under the proximal nail fold. Trauma to the nail complex can cause bleeding under the nail (subungual hematoma), tearing of the nail bed, or separation of the nail plate from the nail bed. Injury to the nail matrix can lead to permanent nail deformity.

\section*{Causes}
\begin{itemize}
\item Crush injury: closing a door on a finger, dropping a heavy object
\item Blunt trauma: hitting the finger with a hammer
\item Catching or tearing the nail on an object
\item Sports injury: jamming a finger
\item Stubbing a toe
\item Repetitive trauma: running or hiking (subungual hematoma)
\item Fingertip amputation (severe lacerations)
\end{itemize}

\section*{When to See a Doctor}
\begin{itemize}
\item Pain and throbbing at the nail site
\item Visible blood under the nail (subungual hematoma)
\item Partial or complete detachment of the nail
\item Laceration or irregularity of the nail bed
\item Swelling and tenderness of the fingertip
\item Blood accumulating under the nail causing pressure
\item Fracture of the distal phalanx
\end{itemize}

\section*{Related Symptoms}
\begin{itemize}
\item Subungual hematoma (blood under nail)
\item Nail avulsion
\item Fingertip injury
\item Finger fracture
\item Nail bed laceration
\item Paronychia (infection after injury)
\end{itemize}

\section*{Self-Care / Home Management}
\begin{itemize}
\item For subungual hematoma: apply ice and elevate
\item Wash minor cuts with soap and water
\item Apply antibiotic ointment and a bandage
\item Take OTC pain relievers
\item Keep the nail clean and dry
\item Do not attempt to drain the hematoma at home
\item Protect the nail from further injury
\end{itemize}

\section*{How Your Doctor Will Evaluate You}
\begin{itemize}
\item Clinical examination of the nail and fingertip
\item X-ray to rule out distal phalanx fracture
\item Assessment of nail bed involvement for repair decisions
\end{itemize}

\section*{Treatment Options}
\begin{itemize}
\item Subungual hematoma: trephination (drainage through the nail)
\item Nail bed laceration: repair with fine absorbable sutures
\item Nail avulsion: replace the nail under the nail fold
\item Splint the fingertip if associated fracture
\item Tetanus prophylaxis if needed
\item Antibiotics for open fractures or dirty wounds
\end{itemize}

\section*{References}
\begin{thebibliography}{99}
\bibitem{ref1} Patel L. "Management of simple nail bed lacerations and subungual hematomas in the emergency department." \textit{Pediatric Emergency Care}, 2014;30(10):742--745.
\bibitem{ref2} Brown RE. "Acute nail bed injuries." \textit{Hand Clinics}, 2002;18(4):561--575.
\bibitem{ref3} Roser SE, Gellman H. "Comparison of nail bed repair versus nail trephination for subungual hematomas in children." \textit{Journal of Hand Surgery}, 1999;24(6):1166--1170.
\bibitem{ref4} Giddins GEB, Hill SJ. "Nail bed injuries." \textit{Trauma}, 2013;15(4):321--330.
\bibitem{ref5} Zook EG, Guy RJ, Russell RC. "A study of nail bed injuries: causes, treatment, and prognosis." \textit{Journal of Hand Surgery}, 1984;9(2):247--252.
\bibitem{ref6} Seaberg DC, Angelos WJ, Paris PM. "Treatment of subungual hematomas with nail trephination: a prospective study." \textit{American Journal of Emergency Medicine}, 1991;9(3):209--210.

\end{thebibliography}
